\section{Introduction}

Acute skeletal muscle injuries are among the most common orthopedic complaints, accounting for a substantial fraction of all musculoskeletal injuries seen in athletic and clinical populations \cite{maffulli2014muscle}. Standard care is largely supportive---rest, ice, compression, elevation, and anti-inflammatory medication---and these interventions rest on a thin mechanistic foundation \cite{deyhle2018role}. Incomplete functional recovery, scar formation, and high recurrence rates remain common, particularly for severe injuries. The underlying difficulty is biological, not simply clinical: healthy muscle regeneration is an emergent, multistage process in which the coordinated interaction of cells and their microenvironment, rather than any single cell type, determines the outcome \cite{forcina2020mechanisms,chazaud2020inflammation}.

The regeneration cascade unfolds over weeks and involves a time-ordered recruitment of neutrophils, monocytes, and macrophages that clear necrotic tissue while releasing signals that govern fibroblast behavior and satellite stem cell (SSC) activation \cite{tidball2017regulation,dumont2015satellite}. As inflammation resolves, macrophages transition from a pro-inflammatory to an anti-inflammatory phenotype, and activated SSCs proliferate, differentiate into myoblasts, and fuse to rebuild damaged fibers \cite{arnold2007inflammatory,bentzinger2013cellular}. Microvascular remodeling occurs in parallel: new capillary sprouts are required to restore perfusion, and angiogenic signaling is tightly coupled to immune cell dynamics \cite{carmeliet2003angiogenesis}. Cytokines---including hepatocyte growth factor (HGF), transforming growth factor beta (TGF-$\beta$), tumor necrosis factor alpha (TNF-$\alpha$), matrix metalloproteinase-9 (MMP-9), vascular endothelial growth factor A (VEGF-A), monocyte chemoattractant protein-1 (MCP-1), and interleukin 10 (IL-10)---form the signaling backbone of this cascade \cite{tatsumi1998hgf,lu2011macrophages}.

Although individual cytokines have been studied extensively, systematically probing the combinatorial space of cytokine manipulations is experimentally intractable. With seven cytokines and realistic variations in decay and diffusion parameters, more than a million distinct combinations are possible. Each experimental arm requires an injury model, timepoint-resolved sampling, and dedicated immunoassays \cite{wehling2009interleukin}, and pleiotropic actions, soluble receptors, and cross-regulation confound interpretation. These barriers motivate an \emph{in silico} approach capable of exploring the combinatorial landscape, generating mechanistic hypotheses, and prioritizing experiments.

In this work, we develop a spatial agent-based model (ABM) of murine skeletal muscle regeneration implemented in the Cellular-Potts framework \cite{graner1992simulation,swat2012compucell3d}. The model represents muscle fibers, satellite stem cells, fibroblasts, neutrophils, macrophages, microvessels, and lymphatic vessels, together with seven diffusing cytokines whose dynamics depend on local collagen density. We calibrate unknown parameters by parameter density estimation using the CaliPro protocol \cite{joslyn2021calipro}, validate held-out outputs against published experimental data, and reproduce eight perturbation experiments drawn from the literature. Finally, we use Latin hypercube sampling combined with partial rank correlation coefficient analysis \cite{marino2008sensitivity} to identify cytokine decay and diffusion parameters whose alteration enhances cross-sectional area (CSA) recovery. The goals are to (i) build a spatial ABM that jointly captures cytokine diffusion and microvascular adaptation, (ii) calibrate and validate against multiple independent datasets, and (iii) use the calibrated model to generate testable hypotheses about cytokine interventions.

\section{Related Work}

Computational modeling of skeletal muscle regeneration has advanced rapidly over the last decade. Early non-spatial ordinary differential equation models captured bulk population dynamics but could not resolve the spatial gradients of cytokines that shape local cell decisions \cite{hillen2009user}. Foundational agent-based models of muscle demonstrated that cell-level rules can reproduce tissue-scale behaviors in healthy muscle and in Duchenne muscular dystrophy \cite{virgilio2015multiscale,martin2017agent}, but these implementations relied on simplified, non-spatial cytokine representations. As a result, prior ABMs could not capture interventions that target cytokine transport (for example, injections of TGF-$\beta$ or recombinant chemokines) nor the interaction between cytokine distributions and local ECM remodeling \cite{urciuolo2013collagen}.

A second limitation of earlier muscle ABMs has been the absence of microvasculature. Capillary perfusion state, angiogenesis, and the balance between VEGF-A and MMP-9 are central to the recovery trajectory \cite{carmeliet2003angiogenesis}, yet most published ABMs of muscle treat the vasculature as a static boundary condition rather than as adaptive agents. The lymphatic system and its role in cytokine clearance have been similarly neglected. Recent ABMs have begun to address cerebral palsy and eccentric contraction injury, but angiogenic remodeling during acute regeneration remains underrepresented.

The Cellular-Potts framework, originally introduced for biological cell sorting \cite{graner1992simulation} and implemented in the CompuCell3D platform \cite{swat2012compucell3d}, is well suited to bridge these gaps because it represents cells as extended lattice domains whose morphology and motility arise from an energy minimization and whose surrounding fields (chemokines, ECM) are solved on the same lattice. Combining this representation with CaliPro-based parameter density estimation \cite{joslyn2021calipro} and LHS/PRCC global sensitivity analysis \cite{marino2008sensitivity} provides a route from literature-derived rules to calibrated, validated, and interpretable predictions. Our model builds on this foundation by (a) introducing explicit cytokine-specific diffusion that responds to local collagen density, (b) adding dynamic capillary and lymphatic agents, and (c) using sensitivity analysis to design \emph{in silico} perturbations that would be infeasible experimentally.
